\begin{enumerate}
    \item Donner la configuration électronique du carbone.
          \textbf{Données~:} $Z(\ce{C})=6$.~\textbf{(0,5pt)}
    \item Combien de liaisons covalentes le carbone doit-il former pour
          respecter la règle de l'octet~?~\textbf{(0,5pt)}
    \item Donner la formule développée plane pour la molécule
          \ce{C2H6O}.~\textbf{(0,5pt)}
    \item Compléter le tableau suivant~:~\textbf{(5,5pt)}
\begin{tcbitemize}[raster columns=3,raster rows=8,raster equal height=rows,
        raster column skip=-0.8pt,raster row skip=-0.8pt,
        raster force size=false,
        raster every box/.style={halign=center,valign=center},
        raster row 1/.style={fontupper=\bfseries},
        raster column 1/.style={add to width=-0.5cm},
        raster column 2/.style={add to width=-1cm},
        raster column 3/.style={add to width=1.5cm},
        ]
    %--------------------------------------------------------------------------
    \tcbitem Non de composé
    \tcbitem Formule brute
    \tcbitem Formule semi-développée
    %--------------------------------------------------------------------------
    \tcbitem
    \tcbitem
    \tcbitem \chemfig{CH_{3}-CH_{3}}
    %--------------------------------------------------------------------------
    \tcbitem
    \tcbitem
    \tcbitem \chemfig{CH_{3}-CH_{2}-CH_{3}}
    %--------------------------------------------------------------------------
    \tcbitem hexane
    \tcbitem
    \tcbitem
    %--------------------------------------------------------------------------
    \tcbitem
    \tcbitem
    \tcbitem \chemfig{CH_{3}-CH(-[:90]CH_3)-CH_{3}}
    %--------------------------------------------------------------------------
    \tcbitem
    \tcbitem
    \tcbitem \chemfig{CH_{3}-CH_2-CH(-[:90]CH_3)-CH_{3}}
    %--------------------------------------------------------------------------
    \tcbitem
    \tcbitem
    \tcbitem \chemfig{CH_{3}-C(-[:90]CH_3)(-[:270]CH_3)-CH_2-CH_{3}}
    %--------------------------------------------------------------------------
    \tcbitem[height=2cm] 2,3-diméthylhexane
    \tcbitem[height=2cm]
    \tcbitem[height=2cm]
    %--------------------------------------------------------------------------
    \tcbitem
    \tcbitem
    \tcbitem \chemfig{CH_{3}-CH_2-CH(-[:90]CH_3)-CH(-[:90]CH_3)-CH_{3}}
    %--------------------------------------------------------------------------
    \tcbitem
    \tcbitem
    \tcbitem \chemfig{CH_{3}-CH(-[:270]CH_2-[:180]CH_3)-CH_2-CH_{3}}
    %--------------------------------------------------------------------------
    \tcbitem
    \tcbitem
    \tcbitem \chemfig{CH_{3}-CH(-[:270]CH_3)-CH(-[:270]CH_2-CH_3)-CH_{3}}
    %--------------------------------------------------------------------------
    \tcbitem[height=2cm] 3-éthylpentane
    \tcbitem[height=2cm]
    \tcbitem[height=2cm]
\end{tcbitemize}
\end{enumerate}
